\documentclass[10pt]{article}
\usepackage[a4paper,margin=1.8cm]{geometry}
\usepackage{amsmath,amssymb}
\usepackage{algorithm}
\usepackage{algpseudocode}
\pagestyle{empty}
\begin{document}

\begin{algorithm}[H]
\caption{f-DAgger}
\label{alg:f-dagger}
\small
\begin{algorithmic}[1]
\Require Expert $\pi_E$, clean data $D_{\mathrm{clean}}$, DAgger rounds $K$
\Ensure Learned policy $\pi^*$
\Statex \textit{Stage 1: collect expert data with piecewise OU noise}
\State $D_{\mathrm{OU}}\gets\emptyset$
\For{each expert episode}
    \State Reset environment; $z\gets0$; $T_A,T_R\gets\infty$; $B\gets\emptyset$
    \State Sample one episode scale $\alpha\sim\mathcal U(0.25,1)$
    \For{each control step $t$}
        \State Observe $o_t$ and query the expert: $a_t^E\gets\pi_E(o_t)$
        \State Update the first gate-activation time $T_A$ and first task-reach time $T_R$
        \State $z\gets\rho z+\sqrt{1-\rho^2}\,(\alpha\sigma\odot\xi)$, where $\xi\sim\mathcal N(0,I)$
        \State Apply the piecewise rule below to obtain $n_t$ from $z$
        \State Save $(o_t,a_t^E)$ in $B$ \Comment{clean expert label}
        \State Execute $\operatorname{clip}(\operatorname{clip}(a_t^E)+n_t)$
    \EndFor
    \State Add $B$ to $D_{\mathrm{OU}}$ only if its sustained-hold check passes
\EndFor
\State $D\gets D_{\mathrm{clean}}\uplus D_{\mathrm{OU}}$
\State $\pi_0\gets\operatorname{TrainFromScratch}(D)$
\Statex \textit{Stage 2: DAgger data aggregation}
\For{$k=1,\ldots,K$}
    \State Roll out student $\pi_{k-1}$ with no additional OU noise
    \State Label every visited observation $o$ with $\pi_E(o)$ to obtain $D_{\mathrm{new}}$
    \State Retain task failures; discard numerical-failure and empty episodes
    \State $D\gets D\uplus D_{\mathrm{new}}$
    \State $\pi_k\gets\operatorname{TrainFromScratch}(D)$
    \State Evaluate $\pi_k$ on the validation scenarios
\EndFor
\State \Return the best validation policy $\pi^*$ among $\pi_1,\ldots,\pi_K$
\end{algorithmic}
\end{algorithm}

\noindent\textbf{Piecewise OU noise.}
Write $z=[z^{\mathrm{arm}};z^{\mathrm{wheel}}]$. The applied perturbation is
\[
n_t=\begin{cases}
[z^{\mathrm{arm}};z^{\mathrm{wheel}}], & t\Delta t<T_A,\\
[z^{\mathrm{arm}};0], & T_A\le t\Delta t<T_R+1.5\ \mathrm{s},\\
0, & t\Delta t\ge T_R+1.5\ \mathrm{s}.
\end{cases}
\]
$T_A$ is the first reach-gate activation; wheel noise remains off thereafter.
$T_R$ is the first task reach: gate active, true end-effector error below 10 mm,
and elapsed time greater than 8 s. An event that never occurs has time $\infty$.

\medskip
\noindent\textbf{Defaults.}
$\Delta t=0.001$ s, $\tau=0.2$ s, $\rho=e^{-\Delta t/\tau}$;
$\sigma$ contains seven arm scales of $0.75$ N\,m and four wheel scales of
$0.15$ rad/s. All channels use independent Gaussian innovations. $K=3$.

\medskip
\noindent\textbf{Collection and training.}
Episodes end 3 s after the first task reach or at 40 s; numerical failures
interrupt collection. Sustained-hold success requires the complete 3 s
post-reach window and error below 10 mm throughout its final second.
DAgger uses fitting scenarios plus fresh low/near-biased development scenarios;
the student executes alone ($\beta=0$), with command limits applied.
$\uplus$ denotes data concatenation. Feature extraction and action-format
conversion are implicit in the observation--expert-action pairs.
Training uses the implemented BC+J objective; all previous data are retained.
Validation uses sustained-hold success, with ties resolved toward later rounds.

\end{document}
